% !TEX root = ../main.tex
\chapter{Models and Results}
\label{chapter:model_results}

\providecommand{\safeplaceholdergraphic}[2][0.92\linewidth]{%
  \IfFileExists{#2.pdf}{\includegraphics[width=#1]{#2.pdf}}{%
    \IfFileExists{#2.png}{\includegraphics[width=#1]{#2.png}}{%
      \IfFileExists{#2.jpg}{\includegraphics[width=#1]{#2.jpg}}{%
        \fbox{\parbox[c][0.24\textheight][c]{0.84\linewidth}{\centering Placeholder for figure\\[0.4em]\texttt{\detokenize{#2}}}}%
      }%
    }%
  }%
}

\providecommand{\safefallbackgraphic}[3][0.92\linewidth]{%
  \IfFileExists{#2.pdf}{\includegraphics[width=#1]{#2.pdf}}{%
    \IfFileExists{#3.pdf}{\includegraphics[width=#1]{#3.pdf}}{%
      \IfFileExists{#2.png}{\includegraphics[width=#1]{#2.png}}{%
        \IfFileExists{#3.png}{\includegraphics[width=#1]{#3.png}}{%
          \IfFileExists{#2.jpg}{\includegraphics[width=#1]{#2.jpg}}{%
            \IfFileExists{#3.jpg}{\includegraphics[width=#1]{#3.jpg}}{%
              \fbox{\parbox[c][0.24\textheight][c]{0.84\linewidth}{\centering Placeholder for figure\\[0.4em]\texttt{\detokenize{#2}}\\[0.4em]\texttt{\detokenize{#3}}}}%
            }%
          }%
        }%
      }%
    }%
  }%
}

\providecommand{\monomerFigTwoTimePaperPath}{\detokenize{C:/Users/leopo/.vscode/thesis_python/jobs/monomer/monomerfig2/validation/homogeneous/07_095914_monomer/figures/time_all_signals_real_imag_abs_monomer}}
\providecommand{\monomerFigTwoTimeLindbladPath}{\detokenize{C:/Users/leopo/.vscode/thesis_python/jobs/monomer/monomerfig2/validation/homogeneous/07_100219_monomer/figures/time_all_signals_real_imag_abs_monomer}}
\providecommand{\monomerFigTwoFreqPaperPath}{\detokenize{C:/Users/leopo/.vscode/thesis_python/jobs/monomer/monomerfig2/validation/homogeneous/07_095914_monomer/figures/freq_all_signals_real_imag_abs_monomer}}
\providecommand{\monomerFigTwoFreqLindbladPath}{\detokenize{C:/Users/leopo/.vscode/thesis_python/jobs/monomer/monomerfig2/validation/homogeneous/07_100219_monomer/figures/freq_all_signals_real_imag_abs_monomer}}
\providecommand{\monomerFigFourTFortyPaperAbsPath}{\detokenize{C:/Users/leopo/.vscode/thesis_python/jobs/monomer/monomer_fig4/01_184049_monomer_fig4_T040_paper_eqs_LEIDER_UEBERSCHRIEBEN/figures/freq_A_abs_monomer_fig4_T040_paper_eqs}}
\providecommand{\monomerFigFourTFortyLindbladAbsPath}{\detokenize{C:/Users/leopo/.vscode/thesis_python/jobs/monomer/monomer_fig4/01_184049_monomer_fig4_T040_paper_eqs_LEIDER_UEBERSCHRIEBEN/figures/freq_A_abs_monomer_fig4_T040_lindblad}}
\providecommand{\monomerFigThreeTZeroLindbladPath}{\detokenize{C:/Users/leopo/.vscode/thesis_python/jobs/monomer/monomer_fig3/01_162945_monomer_fig3_T000_lindblad/figures/freq_all_signals_real_imag_abs_monomer_fig3_T000_lindblad}}
\providecommand{\monomerFigThreeTTwentyLindbladPath}{\detokenize{C:/Users/leopo/.vscode/thesis_python/jobs/monomer/monomer_fig3/01_162659_monomer_fig3_T020_lindblad/figures/freq_all_signals_real_imag_abs_monomer_fig3_T020_lindblad}}
\providecommand{\monomerFigThreeTFortyLindbladPath}{\detokenize{C:/Users/leopo/.vscode/thesis_python/jobs/monomer/monomer_fig3/02_070401_monomer_fig3_T040_lindblad/figures/freq_all_signals_real_imag_abs_monomer_fig3_T040_lindblad}}
\providecommand{\monomerFigFourTZeroLindbladPath}{\detokenize{C:/Users/leopo/.vscode/thesis_python/jobs/monomer/monomer_fig4/01_184018_monomer_fig4_T000_lindblad/figures/freq_all_signals_real_imag_abs_monomer_fig4_T000_lindblad}}
\providecommand{\monomerFigFourTTwentyLindbladPath}{\detokenize{C:/Users/leopo/.vscode/thesis_python/jobs/monomer/monomer_fig4/01_184037_monomer_fig4_T020_lindblad/figures/freq_all_signals_real_imag_abs_monomer_fig4_T020_lindblad}}
\providecommand{\monomerFigFourTFortyLindbladPath}{\detokenize{C:/Users/leopo/.vscode/thesis_python/jobs/monomer/monomer_fig4/01_184055_monomer_fig4_T040_lindblad/figures/freq_all_signals_real_imag_abs_monomer_fig4_T040_lindblad}}
\providecommand{\monomerFigThreeRedfieldPath}{\detokenize{C:/Users/leopo/.vscode/thesis_python/jobs/monomer/monomer_fig3/02_123454_monomer_fig3_homogeneous_redfield_no_rwa_thermal/figures/freq_all_signals_real_imag_abs_monomer_fig2_homogeneous_redfield_no_rwa_thermal}}
\providecommand{\monomerFigFourRedfieldPath}{\detokenize{C:/Users/leopo/.vscode/thesis_python/jobs/monomer/monomer_fig4/03_143823_monomer_fig4_T040_redfield_no_rwa_thermal/figures/freq_all_signals_real_imag_abs_monomer_fig4_T040_redfield_no_rwa_thermal}}

\providecommand{\dimerFigThreeUncoupledPaperPath}{\detokenize{C:/Users/leopo/.vscode/thesis_python/jobs/dimer/dimer_fig3/01_232832_dimer_fig3a_uncoupled_paper_eqs/figures/freq_all_signals_real_imag_abs_dimer_fig3a_uncoupled_paper_eqs}}
\providecommand{\dimerFigThreeCoupledPaperPath}{\detokenize{C:/Users/leopo/.vscode/thesis_python/jobs/dimer/dimer_fig3/01_232845_dimer_fig3b_coupled_paper_eqs/figures/freq_all_signals_real_imag_abs_dimer_fig3b_coupled_paper_eqs_04}}
\providecommand{\dimerFigFourWaitFortySixPaperPath}{\detokenize{C:/Users/leopo/.vscode/thesis_python/jobs/dimer/dimer_fig4/03_092148_dimer_fig4_T046_paper_eqs/figures/freq_all_signals_real_imag_abs_dimer_fig4_T046_paper_eqs}}
\providecommand{\dimerFigFourWaitThreeTenPaperPath}{\detokenize{C:/Users/leopo/.vscode/thesis_python/jobs/dimer/dimer_fig4/03_092213_dimer_fig4_T310_paper_eqs/figures/freq_all_signals_real_imag_abs_dimer_fig4_T310_paper_eqs}}
\providecommand{\dimerFigSixWaitSixtyTwoPaperPath}{\detokenize{C:/Users/leopo/.vscode/thesis_python/jobs/dimer/dimer_fig6/02_172531_dimer_fig6_T062_paper_eqs/figures/freq_all_signals_real_imag_abs_dimer_fig6_T062_paper_eqs}}
\providecommand{\dimerFigThreeUncoupledLindbladPath}{\detokenize{figures/generated/dimer_fig3a_uncoupled_lindblad}}
\providecommand{\dimerFigThreeCoupledLindbladPath}{\detokenize{figures/generated/dimer_fig3b_coupled_lindblad}}
\providecommand{\dimerFigFourWaitFortySixLindbladPath}{\detokenize{figures/generated/dimer_fig4_T046_lindblad}}
\providecommand{\dimerFigFourWaitThreeTenLindbladPath}{\detokenize{figures/generated/dimer_fig4_T310_lindblad}}
\providecommand{\dimerFigSixWaitSixtyTwoLindbladPath}{\detokenize{figures/generated/dimer_fig6_T062_lindblad}}
\providecommand{\dimerFigFourWaitFortySixRedfieldRwaPath}{\detokenize{figures/generated/dimer_fig4_T046_redfield_rwa_thermal}}

This chapter combines the open-system framework of~\autoref{chapter:oqs_open_quantum_systems} with the spectroscopy workflow of~\autoref{chapter:spctr_spectroscopy}. Its central aim is to reproduce selected spectra from Papers~I and II inside one consistent QuTiP-based workflow~\cite{mancaletal2006twodimensionalopticalthreepulse,pisliakovetal2006twodimensionalopticalthreepulse}. In the main text, the exact paper equations runs are kept only for one explicit calibration step per system. After that point, the results are organised in terms of Lindblad reproductions and Redfield extensions. The archived paper equations figure sets are moved to~\autoref{app:paper_eqs_reference_atlas}.

\section{Monomer}
\label{sec:results_monomer}

The discussion starts with the monomer because it is the simplest system in which the complete workflow can be inspected.

\subsection{Model and equations of motion}
\label{subsec:results_monomer_model_eom}

In the strict two-level limit, the third-order signal contains only ground-state bleach and stimulated-emission contributions; excited-state absorption is absent because no higher excited manifold is available, as discussed in~\autoref{sec:spctr_two_level_system_linear_response} and~\autoref{sec:spctr_two_level_third_order_pathways}.

The monomer is the simplest case in which the full three-pulse workflow can be inspected. Its field-free Hamiltonian is, cf.~\cref{eq:spctr_tls_hamiltonian,eq:spctr_tls_dipole},
\begin{equation*}
\hat{H}_{0}^{(\mathrm{mon})}=\omega_{eg}\ket{e}\!\bra{e},
\end{equation*}
and the dipole operator is
\begin{equation*}
\hat{\mu}^{(\mathrm{mon})}=\mu_{ge}\ket{g}\!\bra{e}+\mu_{eg}\ket{e}\!\bra{g}.
\end{equation*}
The published phenomenology of Paper~I is characterised by a population decay rate \(\gamma\) and a pure dephasing rate \(\gamma_{\varphi}\), with \(T_1=\gamma^{-1}\) and \(T_2^{-1}=\gamma/2+\gamma_{\varphi}\). In the QuTiP implementation used here, the same phenomenology is written in Lindblad form with
\begin{equation*}
\hat{L}_{\varphi}=\sqrt{2\gamma_{\varphi}}\ket{e}\!\bra{e},
\qquad
\hat{L}_{\downarrow}=\sqrt{\gamma_{\downarrow}}\ket{g}\!\bra{e},
\qquad
\hat{L}_{\uparrow}=\sqrt{\gamma_{\uparrow}}\ket{e}\!\bra{g}.
\end{equation*}

With the ansatz \(\rho_{eg}=\sigma_{eg}\,\mathrm{e}^{-\mathrm{i}\omega_L t}\), the density-matrix elements obey
\begin{align}
\dot{\sigma}_{eg} &= -\bigl(\Gamma + i\Delta_{L}\bigr)\sigma_{eg}
+ i\,\mu_{eg} E(t)\,\bigl(\rho_{gg}-\rho_{ee}\bigr),
\label{eq:results_monomer_eom_sigma}\\
\dot{\rho}_{ee} &= -\gamma\,\rho_{ee}
+ i\,\mu_{eg}E(t)\sigma_{ge}
- i\,\mu_{eg}E^{*}(t)\sigma_{eg},
\label{eq:results_monomer_eom_ee}\\
\dot{\rho}_{gg} &= - i\,\mu_{eg}E(t)\sigma_{ge}
+ i\,\mu_{eg}E^{*}(t)\sigma_{eg},
\label{eq:results_monomer_eom_gg}
\end{align}
with detuning
\begin{equation*}
\Delta_{L}=\omega_{eg}-\omega_L.
\end{equation*}

For real samples, the measured response reflects an ensemble of local environments rather than a single perfectly identical monomer, so the spectra are evaluated over many static-disorder realisations, see~\autoref{spctr:sec:macroscopic_polarisation}. The transition frequency is sampled from the Gaussian distribution introduced in~\autoref{eq:gaussian_broadening}. The spectra are then obtained from the disorder-averaged signal defined in~\autoref{eq:averaged_response}. The parameter \(\Delta_{\mathrm{inh}}\) denotes the corresponding full width at half maximum of the distribution. Two situations are considered. In the homogeneous case, \(\Delta_{\mathrm{inh}}=0\), all realisations share the same transition frequency and the linewidth is determined by dynamical dephasing and relaxation.

All monomer runs use \(\tilde{\nu}_{eg}=\SI{16000}{\per\centi\metre}\), \(\mu_{ge}=1\), resonant excitation, and pulse amplitudes \((0.005,0.005,0.005)\). The field amplitude is kept sufficiently weak that the excited-state population remains below \(1\%\), ensuring that higher-order nonlinear contributions remain negligible. In units with \(\hbar=1\), the quantity \(\mu_{eg}E(t)\) has units of inverse time and corresponds to the instantaneous Rabi frequency, while the choice \(\mu_{ge}=1\) fixes the field-amplitude scale used throughout this section. For the phenomenological benchmark, the fitted rates are \(\gamma_{\varphi}=0.01~\si{\per\femto\second}\), \(\gamma_{\downarrow}=1/300~\si{\per\femto\second}\), and \(\gamma_{\uparrow}=0\). Static disorder, when included, uses \(\Delta_{\mathrm{inh}}=\SI{200}{\per\centi\metre}\). The main-text role of paper equations is only to verify once that this Lindblad choice reproduces the Paper-I behaviour.

\subsection{Calibration against Paper~I}
\label{subsec:results_monomer_calibration}

\begin{figure}[htbp]
    \centering
    \begin{minipage}[t]{0.49\linewidth}
        \centering
        \safeplaceholdergraphic[0.98\linewidth]{\monomerFigTwoTimePaperPath}
    \end{minipage}\hfill
    \begin{minipage}[t]{0.49\linewidth}
        \centering
        \safeplaceholdergraphic[0.98\linewidth]{\monomerFigTwoTimeLindbladPath}
    \end{minipage}
    \caption{Monomer time-domain calibration at fixed \(t_{\mathrm{coh}}=\SI{300}{\femto\second}\), \(t_{\mathrm{wait}}=\SI{1000}{\femto\second}\), \(N_{\mathrm{inhom}}=1\), and \(\Delta_{\mathrm{inh}}=0\). Left: archived paper equations reference. Right: Lindblad propagation with the fitted rates given above.}
    \label{fig:results_monomer_1d_time_calibration}
\end{figure}

\begin{figure}[htbp]
    \centering
    \begin{minipage}[t]{0.49\linewidth}
        \centering
        \safeplaceholdergraphic[0.98\linewidth]{\monomerFigTwoFreqPaperPath}
    \end{minipage}\hfill
    \begin{minipage}[t]{0.49\linewidth}
        \centering
        \safeplaceholdergraphic[0.98\linewidth]{\monomerFigTwoFreqLindbladPath}
    \end{minipage}
    \caption{Monomer 1D frequency-domain calibration for the same homogeneous setting as in~\cref{fig:results_monomer_1d_time_calibration}. The component structure of the Lindblad run follows the Paper-I-style reference while being generated within the QuTiP workflow.}
    \label{fig:results_monomer_1d_freq_calibration}
\end{figure}

\begin{figure}[htbp]
    \centering
    \begin{minipage}[t]{0.49\linewidth}
        \centering
        \safeplaceholdergraphic[0.98\linewidth]{\monomerFigFourTFortyPaperAbsPath}
    \end{minipage}\hfill
    \begin{minipage}[t]{0.49\linewidth}
        \centering
        \safeplaceholdergraphic[0.98\linewidth]{\monomerFigFourTFortyLindbladAbsPath}
    \end{minipage}
    \caption{Representative 2D monomer comparison at \(T=\SI{40}{\femto\second}\). Only the absorptive spectrum is shown, because it is the cleanest comparison of the diagonally elongated inhomogeneous line shape. The archived full paper equations atlas is moved to~\autoref{app:paper_eqs_reference_atlas}.}
    \label{fig:results_monomer_2d_calibration}
\end{figure}

These three comparisons fix the thesis logic for the monomer. Once the phenomenological rates are matched, the main text no longer needs paper equations. The paper-equation runs are retained only as archived comparison material, while the actual monomer results are presented as Lindblad reproductions and Redfield extensions.

\subsection{Lindblad monomer results}
\label{subsec:results_monomer_lindblad}

\begin{figure}[p]
    \centering
    \safeplaceholdergraphic[0.95\linewidth]{\monomerFigThreeTZeroLindbladPath}
    \vspace{0.8em}
    \safeplaceholdergraphic[0.95\linewidth]{\monomerFigThreeTTwentyLindbladPath}
    \vspace{0.8em}
    \safeplaceholdergraphic[0.95\linewidth]{\monomerFigThreeTFortyLindbladPath}
    \caption{Homogeneous monomer waiting-time series obtained with the calibrated Lindblad model at \(T=\SI{0}{\femto\second}\), \SI{20}{\femto\second}, and \SI{40}{\femto\second}. The common parameters are \(\Delta_{\mathrm{inh}}=0\), \(N_{\mathrm{inhom}}=1\), \(\gamma_{\varphi}=0.01~\si{\per\femto\second}\), \(\gamma_{\downarrow}=1/300~\si{\per\femto\second}\), and pulse amplitudes \((0.005,0.005,0.005)\).}
    \label{fig:results_monomer_2d_hom_lindblad}
\end{figure}

\begin{figure}[p]
    \centering
    \safeplaceholdergraphic[0.95\linewidth]{\monomerFigFourTZeroLindbladPath}
    \vspace{0.8em}
    \safeplaceholdergraphic[0.95\linewidth]{\monomerFigFourTTwentyLindbladPath}
    \vspace{0.8em}
    \safeplaceholdergraphic[0.95\linewidth]{\monomerFigFourTFortyLindbladPath}
    \caption{Inhomogeneous monomer waiting-time series obtained with the same Lindblad model at \(T=\SI{0}{\femto\second}\), \SI{20}{\femto\second}, and \SI{40}{\femto\second}. Here \(\Delta_{\mathrm{inh}}=\SI{200}{\per\centi\metre}\).}
    \label{fig:results_monomer_2d_inhom_lindblad}
\end{figure}

The homogeneous Lindblad series reproduces the compact axis-aligned monomer reference profile, whereas the inhomogeneous series develops the expected diagonal elongation. This is the key Paper-I distinction, but in the present chapter it is attributed to the QuTiP Lindblad implementation rather than to the paper-specific equations of motion. The archived paper equations monomer figures remain available in~\autoref{app:paper_eqs_reference_atlas} for direct comparison.

\subsection{Microscopic monomer extension}
\label{subsec:results_monomer_redfield}

\begin{figure}[htbp]
    \centering
    \begin{minipage}[t]{0.49\linewidth}
        \centering
        \safeplaceholdergraphic[0.98\linewidth]{\monomerFigThreeRedfieldPath}
    \end{minipage}\hfill
    \begin{minipage}[t]{0.49\linewidth}
        \centering
        \safeplaceholdergraphic[0.98\linewidth]{\monomerFigFourRedfieldPath}
    \end{minipage}
    \caption{Representative monomer Redfield calculations that go beyond the phenomenological Paper-I benchmark. The Ohmic bath parameters are \(T/\omega_0=0.0033443363446838082\) (about \SI{77}{\kelvin}), \(\omega_c/\omega_0=100\), \(\alpha/\omega_0=0.0267\), and \(s=1\).}
    \label{fig:results_monomer_redfield_examples}
\end{figure}

These Redfield panels are not meant as replacements for the calibrated Lindblad baseline. They play a different role: they show that the same signal-construction workflow also works with an explicit bath model and therefore goes beyond the phenomenology of Paper~I.

\subsection{Monomer summary}
\label{subsec:results_monomer_summary}

The monomer section first calibrates Lindblad dynamics once against the archived Paper-I-style reference and then keeps only Lindblad and Redfield in the main story. Lindblad is used as the thesis baseline for reproducing the published line-shape logic; Redfield is used to extend the workflow to a microscopic bath model.

\section{Dimer}
\label{sec:results_dimer}

\subsection{Model and calibration strategy}
\label{subsec:results_dimer_model_eom}

The dimer is the minimal model in which cross peaks, waiting-time beats, and adjacent-manifold transitions appear together. In the truncated manifold \(\{\ket{g},\ket{1},\ket{2},\ket{f}\}\), the field-free Hamiltonian is
\begin{equation*}
\hat{H}_{0}^{(\mathrm{dim})}
=
\omega_{1}\ket{1}\!\bra{1}
+
\omega_{2}\ket{2}\!\bra{2}
+
J\bigl(\ket{1}\!\bra{2}+\ket{2}\!\bra{1}\bigr)
+
(\omega_{1}+\omega_{2})\ket{f}\!\bra{f},
\end{equation*}
and the dipole operator connects adjacent manifolds. The main phenomenological dimer parameters follow Paper~II~\cite{pisliakovetal2006twodimensionalopticalthreepulse}: a \(T=0\) uncoupled-versus-coupled comparison and a waiting-time series for the coupled dimer with pulse amplitudes again fixed to \((0.005,0.005,0.005)\).

As in the monomer case, the thesis logic is to use paper equations only once as a calibration reference. The Lindblad rates used to mimic the paper behaviour are the same phenomenological values employed during the dimer tests and notebook checks, \(\gamma_{\varphi}=0.01~\si{\per\femto\second}\), \(\gamma_{\downarrow}=1/300~\si{\per\femto\second}\), and \(\gamma_{\uparrow}=0\). The archived paper-equation panels are collected in~\autoref{app:paper_eqs_reference_atlas}.

\subsection{Calibration against Paper~II}
\label{subsec:results_dimer_calibration}

\begin{figure}[htbp]
    \centering
    \begin{minipage}[t]{0.49\linewidth}
        \centering
        \safeplaceholdergraphic[0.98\linewidth]{\dimerFigFourWaitFortySixPaperPath}
    \end{minipage}\hfill
    \begin{minipage}[t]{0.49\linewidth}
        \centering
        \safeplaceholdergraphic[0.98\linewidth]{\dimerFigFourWaitFortySixLindbladPath}
    \end{minipage}
    \caption{Single explicit dimer calibration step at \(T=\SI{46}{\femto\second}\). Left: archived paper equations reference. Right: the corresponding Lindblad comparison that calibrates the phenomenological rates inside the QuTiP workflow.}
    \label{fig:results_dimer_calibration}
\end{figure}

This subsection plays the same role as the monomer calibration: it makes clear that the thesis aims to reproduce the Paper-II behaviour inside QuTiP and then switch to Lindblad and Redfield as the only main-text solvers.

\subsection{Lindblad dimer results}
\label{subsec:results_dimer_lindblad}

\begin{figure}[p]
    \centering
    \begin{minipage}[t]{0.49\linewidth}
        \centering
        \safeplaceholdergraphic[0.98\linewidth]{\dimerFigThreeUncoupledLindbladPath}
    \end{minipage}\hfill
    \begin{minipage}[t]{0.49\linewidth}
        \centering
        \safeplaceholdergraphic[0.98\linewidth]{\dimerFigThreeCoupledLindbladPath}
    \end{minipage}
    \caption{Lindblad reproduction of the Paper-II \(T=0\) uncoupled-versus-coupled comparison. The intended interpretation is the same as in the archived paper-equation reference: no strong cross peaks in the uncoupled control and clear cross peaks already at \(T=0\) for the coupled dimer.}
    \label{fig:results_dimer_lindblad_fig3}
\end{figure}

\begin{figure}[p]
    \centering
    \safeplaceholdergraphic[0.95\linewidth]{\dimerFigFourWaitFortySixLindbladPath}
    \vspace{0.8em}
    \safeplaceholdergraphic[0.95\linewidth]{\dimerFigSixWaitSixtyTwoLindbladPath}
    \vspace{0.8em}
    \safeplaceholdergraphic[0.95\linewidth]{\dimerFigFourWaitThreeTenLindbladPath}
    \caption{Representative Lindblad waiting-time panels for the coupled dimer: inhomogeneous \(T=\SI{46}{\femto\second}\), homogeneous companion \(T=\SI{62}{\femto\second}\), and long-time inhomogeneous \(T=\SI{310}{\femto\second}\).}
    \label{fig:results_dimer_lindblad_waiting}
\end{figure}

In this way the main dimer story is meant to parallel the paper logic without remaining tied to the paper-specific solver: the uncoupled control isolates the diagonal peaks, the coupled \(T=0\) panel fills the cross-peak positions, and the waiting-time series separates the coherence-dominated from the relaxation-dominated regime.

\subsection{Microscopic Redfield extension}
\label{subsec:results_dimer_redfield}

\begin{figure}[htbp]
    \centering
    \safeplaceholdergraphic[0.95\linewidth]{\dimerFigFourWaitFortySixRedfieldRwaPath}
    \caption{Representative microscopic dimer Redfield calculation in the rotating-wave setting. The thesis-facing bath parameters are \(T/\omega_0=0.0033443363446838082\) (about \SI{77}{\kelvin}), \(\omega_c/\omega_0=100\), \(\alpha/\omega_0=0.0267\), and \(s=1\), with the secular cutoff reported in normalized form as \(\texttt{sec\_cutoff}/\omega_0\). The intended main-text choice is a very small cutoff, \(\texttt{sec\_cutoff}/\omega_0 \approx 10^{-5}\), together with \texttt{rwa\_sl=true}.}
    \label{fig:results_dimer_redfield_rwa}
\end{figure}

The Redfield dimer result is not used as another phenomenological fit. Its purpose is to document the next step beyond the paper reproduction: a direct bath-based propagation inside the QuTiP implementation, with the secular cutoff stated in the same normalized units as the bath parameters.

\subsection{Dimer summary}
\label{subsec:results_dimer_summary}

The dimer section follows the same logic as the monomer but with a richer spectrum: one explicit paper equations comparison calibrates the phenomenological Lindblad baseline, the intended main-text reproduction is then Lindblad-only, and the microscopic extension is formulated in terms of the rotating-wave Redfield implementation with a normalized secular cutoff.

\section{Towards the general \texorpdfstring{$N$}{N}-site extension}
\label{sec:results_nsite_extension}

The numerical workflow developed in this chapter is not restricted to the monomer and dimer. At the level of the implementation described in~\autoref{app:numerical_implementation}, the same sequence remains available for a general \(N\)-site excitonic model: construct the Hamiltonian and dipole operators in the truncated excitation manifolds, propagate under the three-pulse field, isolate the desired phase-matched contribution by phase cycling, average over static disorder, and Fourier transform the resulting signal into \((\omega_{\mathrm{coh}},\omega_{\mathrm{det}})\) space.

What remains unfinished in the present thesis is not the basic simulation machinery, but the final production of a larger, fully interpreted figure set. A trimer remains the natural next benchmark because it would test the same pipeline in the first system where multiple cross-peak families compete simultaneously while the spectrum is still interpretable by hand.

\section{Chapter summary}
\label{sec:results_summary}

This chapter is now organised around one consistent rule: use paper equations only to calibrate, then present the thesis results in terms of Lindblad reproductions and Redfield extensions. For the monomer, this structure is already fully visible with existing data. For the dimer, the same structure is now built into the thesis text and figure slots, while the archived paper-equation panels remain available in~\autoref{app:paper_eqs_reference_atlas} so that the reader can always compare the new QuTiP-based results against the phenomenological reference atlas.

